\textbf{Data and task.} The PhysioNet/Computing in Cardiology Challenge 2012 released 48-hour ICU records with in-hospital mortality labels through PhysioNet~\cite{goldberger2000physiobank}; it has become a staple benchmark for irregular time-series models. We use only the open training set-a.

\textbf{Missing values in clinical series.} Lipton et al.~\cite{lipton2016modeling} showed that missingness indicators alone carry predictive signal in ICU data, because clinicians order measurements in response to patient state. GRU-D~\cite{che2016recurrent} operationalises this with masks, time-deltas and learned decay, and reported gains over GRU baselines with simple imputation and over feature-based baselines on PhysioNet 2012 and MIMIC-III. Later work replaces the recurrent core with interpolation networks~\cite{shukla2019interpolation}, set functions over observations~\cite{horn2019set}, latent ODEs~\cite{rubanova2019latent}, graph-guided message passing~\cite{zhang2021graph} and self-supervised transformers~\cite{tipirneni2021self}. We do not reimplement these; GRU-D is the standard reference they are compared with.

\textbf{Benchmarks and strong tabular baselines.} The MIMIC-III benchmark~\cite{harutyunyan2017multitask}, HiRID-ICU-Benchmark~\cite{yche2021hirid} and the YAIB framework~\cite{water2023yet} include logistic regression or gradient-boosted trees on engineered features and find them competitive with recurrent and attention models. Grinsztajn et al.~\cite{grinsztajn2022why} analyse why tree ensembles remain strong on tabular data of medium size. Gradient boosting itself follows~\cite{chen2016xgboost}; we use scikit-learn's histogram implementation.

\textbf{What is left open.} We do not survey how the architecture papers cited above tune their baselines. Our study is a small controlled comparison of a tuned gradient-boosting baseline, a fixed tuning protocol across systems, and a decomposition of the contribution of mask, delta and decay, on one cohort and with several seeds; it cannot replace multi-site validation.
